\documentclass[11pt]{article}
\usepackage[margin=0.9in]{geometry}
\usepackage{amsmath,amssymb,amsthm,mathtools}
\usepackage{booktabs}
\usepackage{longtable}
\usepackage{enumitem}
\usepackage[hidelinks]{hyperref}
\usepackage{microtype}

% This paper carries roughly thirty labelled environments in ten pages; the
% default inter-theorem skip is reduced so that the epistemic labelling does not
% by itself inflate the page count.
\makeatletter
\def\thm@space@setup{\thm@preskip=4pt \thm@postskip=4pt}
\makeatother
\theoremstyle{plain}
\setlist{itemsep=1.5pt,topsep=3pt,parsep=0pt}

\newtheorem{definition}{Definition}[section]
\newtheorem{axiom}[definition]{Axiom}
\newtheorem{assumption}[definition]{Assumption}
\newtheorem{conjecture}[definition]{Conjecture}
\newtheorem{lemma}[definition]{Lemma}
\newtheorem{proposition}[definition]{Proposition}
\newtheorem{theorem}[definition]{Theorem}
\newtheorem{corollary}[definition]{Corollary}
\newtheorem{interpretation}[definition]{Interpretation}
\newtheorem{empirical}[definition]{Empirical statement}
\newtheorem{established}[definition]{Established physical result}
\newtheorem{newresult}[definition]{New result claimed by this work}
\newtheorem{counterexample}[definition]{Counterexample}
\newtheorem{failuremode}[definition]{Failure mode}
\newtheorem{disagreement}[definition]{Preserved disagreement}

\newcommand{\R}{\mathbb{R}}
\newcommand{\Lie}{\mathcal{L}}
\newcommand{\pr}{\pi}

\title{Layered Composition, Reciprocity, and Protocol-Relative Realization:\\
a Synthesis of the PRINCIPIA PHYSICA Topic Series}
\author{PRINCIPIA PHYSICA Synthesis Worker}
\date{15 August 2026}

\begin{document}
\maketitle

\begin{abstract}
Three topic papers of this pilot treat compositional structure, Hamiltonian
reconstruction with interconnection, and empirical realization with its limits.
This paper synthesizes them. We first fix what a synthesis may claim: it may
translate, compare, and combine typed claims, but it may not upgrade the
epistemic type of an inherited claim, and agreement among sources is never
evidence. We then define one layered theory object into which the three
vocabularies translate, and prove three results visible only once the papers are
placed side by side. First, for the declared linear interconnection the
composites split into exactly three classes relative to the product symplectic
form: decoupled, reciprocally coupled (Hamiltonian but not factorized), and
non-reciprocal (not Hamiltonian for that form); the middle class is nonempty, so
``composable'' and ``factorizing'' differ. Second, isomorphism of closed systems
does not descend to the protocol quotient: two symplectomorphic Hamiltonian
systems can have different declared port readings, so assigning
protocol-equivalence classes to isomorphism classes is not well defined. Third,
limits need not commute with composition: the $\omega\to0$ limit of a coupled
oscillator pair carries an interaction no limit of the isolated factors
supplies, and the discrepancy is exponential. We record what the accompanying
Lean~4 and Haskell drafts do and do not establish, and preserve the source
series' disagreements and unverified items rather than resolving them by fiat.
\end{abstract}

\section{What a synthesis is permitted to claim}
\label{sec:discipline}

This document combines three previously issued topic papers of the PRINCIPIA
PHYSICA pilot: \emph{Compositional Systems in Finite-Dimensional Classical
Mechanics} (cited below as CS), \emph{Hamiltonian Reconstruction, Reciprocity,
and Typed Physical Structure} (HR), and \emph{Empirical Realization,
Equivalence, Limits, and Falsification} (EL). Claim identifiers are prefixed
accordingly; identifiers prefixed S belong to this paper.

\begin{assumption}[Inherited scope, S-A1]
\label{ass:scope}
All state spaces are finite-dimensional, Hausdorff, second-countable smooth
manifolds; all maps carry the differentiability stated for them; flows are
asserted only where their existence is stated or proved. Field theory,
infinite-dimensional geometry, general relativity, and quantum theory are
outside the proved scope and appear only as boundary remarks. This assumption is
inherited unchanged from CS-A1, HR-A2, and EL-A1 and is not defended here.
\end{assumption}

\begin{axiom}[Synthesis discipline, S-A2]
\label{ax:discipline}
A synthesis may (i) translate a claim between source vocabularies, (ii) combine
claims whose hypotheses are jointly satisfiable, and (iii) prove new claims from
inherited ones. It may not (iv) raise the epistemic type of an inherited claim,
(v) treat agreement between sources as evidence for either, or (vi) discard a
source disagreement, counterexample, or failed construction. Items (iv)--(vi)
are stipulations governing this document, not assertions about the world.
\end{axiom}

Axiom~\ref{ax:discipline} has an immediate consequence for how the pilot's own
review record may be used.

\begin{empirical}[Process record and its evidential weight, S-E1]
\label{emp:review}
The supplied review artifact \texttt{topic-series-review.json} records verdict
\texttt{ACCEPT}, confidence $1$, an empty list of blocking findings, an empty
list of counterexamples, and an empty list of required changes, with thirteen
per-claim assessments all marked \texttt{supported}. Every one of those thirteen
assessments carries the identical explanation string, ``Claim aligns with
explicit Epistemic rules.'' This records what the review process emitted. By
S-A2(v) it is not evidence that any assessed claim is true: an explanation
string that does not mention the content of the claim it assesses distinguishes
no claim from any other, and confidence $1$ with zero counterexamples is
consistent with a review that did not attempt refutation. The statuses in
Table~\ref{tab:inherited} were therefore obtained by re-derivation here.
\end{empirical}

\begin{interpretation}[Founding thesis, S-I1]
\label{int:thesis}
The sentence ``physics is the study of empirically realized compositional
mathematical structures'' is treated throughout as a methodological proposal
under test. This paper tests one consequence of it: whether ``compositional''
and ``empirically realized'' can be kept apart under pressure. Results S-N1,
S-N2, and S-N3 are all separations, and each is therefore evidence that the two
predicates are not interchangeable. None of them establishes that the thesis is
correct, exhaustive, or unique.
\end{interpretation}

\section{Inventory of inherited claims}
\label{sec:inventory}

Table~\ref{tab:inherited} lists the inherited claims used here and the status
each has \emph{in this round}. ``Re-derived'': the stated computation was
carried out independently here and agreed with the source. ``Inherited'': used
on the source's own authority, not rechecked. ``Convention-dependent'': the
truth value depends on a normalization this paper does not adjudicate.

\begin{table}[h]
\centering
\footnotesize
\begin{tabular}{@{}llp{6.6cm}l@{}}
\toprule
Id & Type & Content (abbreviated) & Status here \\
\midrule
CS-D1,D5 & definition & kinematical system; parallel composite & inherited \\
CS-P1,P3 & proposition & category of typed processes; symmetric monoidal & skeleton machine-checked \\
CS-L1 & lemma & product two-form is symplectic & nondegeneracy half checked \\
CS-T1,C1 & theorem, cor. & additive Hamiltonian generates factorized flow & skeleton machine-checked \\
CS-P4 & proposition & interaction is not determined by the factors & re-derived \\
CS-N1 & new result & four composition layers are logically distinct & inherited \\
HR-P1 & proposition & $d\omega(X_f,X_g,X_h)=2\operatorname{Jac}(f,g,h)$ & convention-dependent \\
HR (c.ex.) & counterex. & Leibniz bracket on $\R^3$ with Jacobiator $x_1$ & re-derived \\
HR-T1 & theorem & symplectic flow is locally Hamiltonian; $H^1_{\mathrm{dR}}$ obstruction & inherited \\
HR-N1 & new result & $\Lie_Y\omega_\oplus=0 \iff C_{21}=C_{12}^{T}$ & re-derived \\
HR-N2,C2 & new result, cor. & split obstruction $a_1c_{12}=a_2c_{21}$ & re-derived \\
HR-N3 & new result & port-sensitive non-conflation of masses & re-derived \\
HR-P3 & proposition & divergence vanishes for all couplings & re-derived \\
EL-D1,D2 & definition & protocol; exact and $\varepsilon$ observational equivalence & inherited \\
EL-P1,N1 & prop., new res. & $\sim_P$ is an equivalence; quotient is maximal & skeleton machine-checked \\
EL (c.ex.) & counterex. & $\varepsilon$-equivalence is not transitive & re-derived \\
EL-P2 & proposition & transported interpretation $\Rightarrow$ $\sim_P$ & skeleton machine-checked \\
EL-N2,C1 & new result, cor. & finite samples underdetermine smooth trajectories & re-derived \\
EL-P3,P4 & proposition & Newtonian and free limits hold only on bounded domains & re-derived \\
EL-T1,C2 & theorem, cor. & $\sim_P$ $\Rightarrow$ equal falsification status & skeleton machine-checked \\
\bottomrule
\end{tabular}
\caption{Inherited claims and their verification status in this synthesis round.
``Skeleton machine-checked'' means the order-and-logic content was formalized and
checked in \texttt{Proofs.lean}; the differential-geometric content was not (see
Section~\ref{sec:formal}).}
\label{tab:inherited}
\end{table}

\section{One layered object}
\label{sec:object}

CS separates kinematical product, observable embedding, dynamical
factorization, and empirical compositionality (CS-N1). HR separates kinematics,
dynamics, interface, and empirical layer (HR-D1). EL separates preparations,
evolutions, measurement kernels, and an adequacy rule (EL-D1). These are three
partitions of overlapping material. The following definition fixes one carrier
for all three.

\begin{definition}[Realized theory object, S-D1]
\label{def:object}
A \emph{realized theory object} is a tuple $\mathcal{R}=(K,D,I,E,P,A)$ in which
\begin{enumerate}[label=(\alph*),leftmargin=2.2em]
\item $K$ is the kinematical layer: a state manifold $X$, a designated
      observable algebra $\mathcal{O}\subseteq C^\infty(X,\R)$, and any
      geometric structure declared on $X$ (here, a symplectic form $\omega$);
\item $D$ is the dynamical layer: a declared generator on $X$ together with the
      domain on which its flow is asserted to exist;
\item $I$ is the interface layer: a finite list of \emph{ports}, that is,
      distinguished elements of $\mathcal{O}$ through which composition with
      other objects and intervention from outside are permitted to act;
\item $E$ is the empirical layer: preparation labels with their represented
      distributions, measurement labels with their kernels, calibration
      conventions fixing dimensions and units, and declared uncertainties;
\item $P$ is a protocol in the sense of EL-D1, indexing admitted
      (preparation, intervention, measurement, time) tuples;
\item $A$ is a predeclared adequacy rule assigning to each protocol index an
      acceptance region determined by the prediction at that index and by the
      declared error model alone.
\end{enumerate}
Composition of realized theory objects is declared componentwise on $K$ and $P$
and by explicit additional data on $D$ and $I$; the content of
Section~\ref{sec:trichotomy} is that this additional data is not optional.
\end{definition}

\begin{proposition}[Vocabulary translation is well defined and non-collapsing, S-P1]
\label{prop:translate}
The assignment $\tau$ of Table~\ref{tab:layers} sends each primitive notion of
CS, HR, and EL to exactly one component of Definition~\ref{def:object}, and
sends the four layers of CS-N1, the four layers of HR-D1, and the four data
items of EL-D1 to pairwise distinct components within each source.
\end{proposition}

\begin{proof}
The claim is a statement about the finite table below, and is verified by
inspecting it. Each row of Table~\ref{tab:layers} names one source notion and
one target component, so $\tau$ is a function. Reading the table by source: the
CS entries occupy $K$, $K$, $D$, and $E{+}P$; the HR entries occupy $K$, $D$,
$I$, and $E$; the EL entries occupy $E$, $D$, $E$, and $A$. Within each source
the listed targets are pairwise distinct, which is the second assertion. No
source notion appears in two rows, which is the first.
\end{proof}

\begin{table}[h]
\centering
\small
\begin{tabular}{@{}lll@{}}
\toprule
Source notion & Source id & Component of S-D1 \\
\midrule
kinematical product & CS-N1 layer 1 & $K$ \\
observable embedding & CS-N1 layer 2 & $K$ (algebra $\mathcal{O}$) \\
dynamical factorization & CS-N1 layer 3 & $D$ \\
empirical compositionality & CS-N1 layer 4 & $E$ with $P$ \\
kinematic layer & HR-D1 $K$ & $K$ \\
dynamical layer & HR-D1 $D$ & $D$ \\
port observables & HR-D1 $I$ & $I$ \\
preparations, dimensions, uncertainties & HR-D1 $E$ & $E$ \\
preparation distribution $\pi_a$ & EL-D1(ii) & $E$ \\
evolutions $\Phi^i_{t,s}$ & EL-D1(i) & $D$ \\
measurement kernel $K_j$ & EL-D1(iii) & $E$ \\
adequacy rule & EL-D1(iv) & $A$ \\
\bottomrule
\end{tabular}
\caption{The translation $\tau$ of Proposition~\ref{prop:translate}.}
\label{tab:layers}
\end{table}

Proposition~\ref{prop:translate} is bookkeeping: it asserts that the three
vocabularies can be carried by one tuple without collision, and nothing more. It
does not assert that the three layer separations are \emph{equivalent}, and it
supplies no functor, no universal property, and no uniqueness statement.

\section{The composition trichotomy}
\label{sec:trichotomy}

Fix two symplectic factors $(\R^{2n_i},\omega_i)$, $i=1,2$, with canonical
coordinates $(q_i,p_i)$ and $\omega_i=\sum_a dq_{ia}\wedge dp_{ia}$. Write
$\omega_\oplus=\pr_1^*\omega_1+\pr_2^*\omega_2$, which is symplectic by CS-L1.
Let $h_0=\pr_1^*h_1+\pr_2^*h_2$ be an additive Hamiltonian and let
\begin{equation}\label{eq:Y}
Y=(0,\;-C_{12}q_2,\;0,\;-C_{21}q_1)
\end{equation}
be the linear interconnection field of HR, with $C_{12}\in\R^{n_1\times n_2}$
and $C_{21}\in\R^{n_2\times n_1}$. The total field is $Z=X_{h_0}+Y$.

\begin{lemma}[Interconnection contraction, S-L1]
\label{lem:contraction}
$\displaystyle \iota_Y\omega_\oplus=\sum_{a}(C_{12}q_2)_a\,dq_{1a}
 +\sum_{b}(C_{21}q_1)_b\,dq_{2b}$, and
\begin{equation}\label{eq:defect}
 d\,\iota_Y\omega_\oplus=\sum_{a,b}\big[(C_{21})_{ba}-(C_{12})_{ab}\big]\,
 dq_{1a}\wedge dq_{2b}.
\end{equation}
\end{lemma}

\begin{proof}
For a single canonical pair, $\iota_Y(dq\wedge dp)=(dq(Y))\,dp-(dp(Y))\,dq$.
By \eqref{eq:Y} the $q$-components of $Y$ vanish and its $p_1$-component is
$-C_{12}q_2$, so the first factor contributes $\sum_a(C_{12}q_2)_a\,dq_{1a}$;
the second factor contributes $\sum_b(C_{21}q_1)_b\,dq_{2b}$ by the same
computation. Applying $d$ and using $(C_{12}q_2)_a=\sum_b(C_{12})_{ab}q_{2b}$,
$(C_{21}q_1)_b=\sum_a(C_{21})_{ba}q_{1a}$ gives
\[
 \sum_{a,b}(C_{12})_{ab}\,dq_{2b}\wedge dq_{1a}
 +\sum_{a,b}(C_{21})_{ba}\,dq_{1a}\wedge dq_{2b},
\]
and antisymmetry of the wedge collects the two sums into \eqref{eq:defect}.
\end{proof}

\begin{newresult}[Composition trichotomy relative to $\omega_\oplus$, S-N1]
\label{res:trichotomy}
Exactly one of the following holds for the pair $(C_{12},C_{21})$, and each case
is realized by an explicit choice of couplings.
\begin{enumerate}[label=(\roman*),leftmargin=2.2em]
\item \emph{Decoupled}: $C_{12}=0$ and $C_{21}=0$. Then $Z=X_{h_0}$ splits as
      $(X_{h_1},X_{h_2})$, and wherever the component flows exist the flow of
      $Z$ is their product.
\item \emph{Reciprocally coupled}: $C_{21}=C_{12}^{T}$ and $C_{12}\neq0$. Then
      $Z$ is Hamiltonian for $\omega_\oplus$, namely $Z=X_{h_0+H_{\rm int}}$ with
      $H_{\rm int}=q_1^{T}C_{12}q_2$, but $Z$ does not split.
\item \emph{Non-reciprocal}: $C_{21}\neq C_{12}^{T}$. Then
      $\Lie_Z\omega_\oplus\neq0$, so $Z$ is not Hamiltonian for $\omega_\oplus$
      and is not even symplectic for it.
\end{enumerate}
Consequently ``the composite is Hamiltonian'' and ``the composite factorizes''
are inequivalent conditions: case (ii) satisfies the first and violates the
second.
\end{newresult}

\begin{proof}
The three cases are exhaustive and mutually exclusive as conditions on
$(C_{12},C_{21})$, since (i) and (ii) are the disjoint subcases $C_{12}=0$ and
$C_{12}\neq0$ of $C_{21}=C_{12}^{T}$ (note that $C_{12}=0$ together with
$C_{21}=C_{12}^{T}$ forces $C_{21}=0$), and (iii) is its negation.

For (i), $Z=X_{h_0}$ and CS-T1 gives $X_{h_0}(x,y)=(X_{h_1}(x),X_{h_2}(y))$;
CS-C1 gives the flow statement on the domain where both component flows exist.

For (ii), $H_{\rm int}=q_1^{T}C_{12}q_2$ has
$dH_{\rm int}=\sum_a(C_{12}q_2)_a\,dq_{1a}+\sum_b(C_{12}^{T}q_1)_b\,dq_{2b}$,
which equals $\iota_Y\omega_\oplus$ of Lemma~\ref{lem:contraction} precisely
when $C_{21}=C_{12}^{T}$; nondegeneracy of $\omega_\oplus$ then gives
$Y=X_{H_{\rm int}}$ and hence $Z=X_{h_0+H_{\rm int}}$. That $Z$ does not split:
the $p_1$-component of $Z$ is $-\partial_{q_1}h_1-C_{12}q_2$, which is an affine
function of $q_2$ with linear part $C_{12}\neq0$ and therefore is not constant
in the second factor's coordinates. A field that splits as $(V_1(x_1),V_2(x_2))$
has every first-factor component independent of $x_2$.

For (iii), $X_{h_0}$ preserves $\omega_\oplus$, so
$\Lie_Z\omega_\oplus=\Lie_Y\omega_\oplus=d\,\iota_Y\omega_\oplus$ by Cartan's
formula and $d\omega_\oplus=0$. The two-forms $dq_{1a}\wedge dq_{2b}$ are
linearly independent, so \eqref{eq:defect} vanishes exactly when
$C_{21}=C_{12}^{T}$, which fails by hypothesis. A field failing to preserve
$\omega_\oplus$ cannot be Hamiltonian for it.

Realization: $n_1=n_2=1$ with $(c_{12},c_{21})$ equal to $(0,0)$, $(3,3)$, and
$(3,0)$ realizes (i), (ii), (iii) respectively.
\end{proof}

\begin{corollary}[Factorization is strictly stronger than Hamiltonicity, S-C1]
\label{cor:strict}
Within the family \eqref{eq:Y}, the composites whose flow factorizes are exactly
those in class (i), and they form a proper subset of the composites that are
Hamiltonian for $\omega_\oplus$, which are exactly classes (i) and (ii).
\end{corollary}

\begin{proof}
Immediate from S-N1: class (i) factorizes; class (ii) is Hamiltonian and does
not factorize; class (iii) is not Hamiltonian. Class (ii) is nonempty, so the
inclusion is proper.
\end{proof}

Corollary~\ref{cor:strict} is the precise sense in which CS-N1's third layer
(dynamical factorization) sits strictly below the condition that the composite
be a Hamiltonian system at all. It also delimits HR-N1: reciprocity buys
Hamiltonicity, not independence.

\begin{counterexample}[The trichotomy is relative to $\omega_\oplus$]
Class (iii) says only that $Z$ is not Hamiltonian for the \emph{product} form.
It leaves open whether $Z$ is Hamiltonian for some other symplectic form. In one
degree of freedom per factor, HR-N2 closes this question for forms respecting
the declared split, and Proposition~\ref{prop:reprove} below re-derives it; for
forms that mix the factors, and in higher dimension, the question is open. Any
reading of S-N1 as ``non-reciprocal composites are not Hamiltonian'' without the
qualifier is therefore stronger than what is proved.
\end{counterexample}

\begin{proposition}[Independent re-derivation of the split obstruction, S-P2]
\label{prop:reprove}
Let $n_1=n_2=1$, let the linearization at a fixed point be
\[
 A=\begin{pmatrix}
 0 & 1/m_1 & 0 & 0\\
 -k_1 & 0 & -c_{12} & 0\\
 0 & 0 & 0 & 1/m_2\\
 -c_{21} & 0 & -k_2 & 0
 \end{pmatrix},
 \qquad
 \Omega_0=\begin{pmatrix} a_1 J & 0\\ 0 & a_2 J\end{pmatrix},
 \quad J=\begin{pmatrix}0&1\\-1&0\end{pmatrix},
\]
with $a_1a_2\neq0$. Then $A^{T}\Omega_0+\Omega_0A=0$ holds if and only if
$a_1c_{12}=a_2c_{21}$. Consequently a split-respecting invariant form exists at
the fixed point if and only if $c_{12}=0$ exactly when $c_{21}=0$; when both are
nonzero, $(a_1,a_2)=(c_{21},c_{12})$ is a witness.
\end{proposition}

\begin{proof}
Since $\Omega_0^{T}=-\Omega_0$, we have
$A^{T}\Omega_0=-A^{T}\Omega_0^{T}=-(\Omega_0A)^{T}$, so the invariance equation
is equivalent to $M:=\Omega_0A$ being symmetric. Multiplying out,
\[
 M=\begin{pmatrix}
 -a_1k_1 & 0 & -a_1c_{12} & 0\\
 0 & -a_1/m_1 & 0 & 0\\
 -a_2c_{21} & 0 & -a_2k_2 & 0\\
 0 & 0 & 0 & -a_2/m_2
 \end{pmatrix}.
\]
All off-diagonal entries of $M$ vanish except $M_{13}=-a_1c_{12}$ and
$M_{31}=-a_2c_{21}$, so symmetry of $M$ is exactly $a_1c_{12}=a_2c_{21}$.
For the consequence: if $c_{12}=0$ then $a_2c_{21}=0$ and $a_2\neq0$ force
$c_{21}=0$, and symmetrically; conversely if both vanish any $a_1,a_2\neq0$
work, and if both are nonzero then $(a_1,a_2)=(c_{21},c_{12})$ satisfies
$c_{21}c_{12}=c_{12}c_{21}$ with both entries nonzero.
\end{proof}

This re-derivation agrees with HR-N2 and HR-C2, and yields one datum not
recorded there: the obstruction is independent of $m_1,m_2,k_1,k_2$, since those
parameters occupy only entries of $M$ that are already symmetric. It is a
property of the coupling pattern alone, not of the tuning of the factors.

\section{Isomorphism does not descend to the protocol quotient}
\label{sec:descent}

EL-N1 shows that the quotient by exact observational equivalence $\sim_P$ is the
finest distinction a protocol can draw; EL-P2 shows that an isomorphism yields
$\sim_P$ \emph{provided} it transports the empirical interpretation; HR-N3 shows
that closed-system isomorphism can forget a port. Together they yield a sharper
negative statement than any one alone.

\begin{newresult}[Non-descent, S-N2]
\label{res:nondescent}
Write $\cong$ for isomorphism of closed Hamiltonian systems, that is, existence
of a symplectomorphism intertwining the Hamiltonians, and write $\sim_P$ for
exact observational equivalence under a fixed protocol $P$ with a fixed
empirical layer. Then $\cong$ does not refine $\sim_P$: there exist realized
theory objects $\mathcal{R}_1,\mathcal{R}_2$ with $K,D$ parts satisfying
$\cong$ and with $\mathcal{R}_1\not\sim_P\mathcal{R}_2$. Consequently there is no
map assigning to each $\cong$-class a $\sim_P$-class; the carrier of protocol
content is the pair (mathematical layer, empirical layer), not the mathematical
layer alone.
\end{newresult}

\begin{proof}
Take $X=\R^2$ with $\omega=dq\wedge dp$, $\omega_0=1$, and
$H_m(q,p)=\tfrac{p^2}{2m}+\tfrac{m\omega_0^2q^2}{2}$, for $m=1$ and $m'=4$. Let
$\alpha=\sqrt{m/m'}=\tfrac12$ and $F(q,p)=(\alpha q,\,p/\alpha)$. Then
$F^*\omega=\omega$ because $\det\begin{psmallmatrix}\alpha&0\\0&\alpha^{-1}
\end{psmallmatrix}=1$, and substituting $q=Q/\alpha$, $p=\alpha P$ gives
\[
 H_m(q,p)=\frac{\alpha^2P^2}{2m}+\frac{m\omega_0^2Q^2}{2\alpha^2}
 =\frac{P^2}{2m'}+\frac{m'\omega_0^2Q^2}{2}=H_{m'}(Q,P),
\]
using $\alpha^2=m/m'$. Hence the closed systems satisfy $\cong$.

Now fix the empirical layer for both objects identically: the preparation sets
the declared coordinates to $(q,p)=(0,1)$, the measurement reads the coordinate
$q$ in the same calibrated units, and the admitted time is $t=1$. The flow of
$H_m$ from $(0,1)$ gives $q_m(t)=\frac{1}{m\omega_0}\sin(\omega_0 t)$, so
$q_1(1)=\sin 1\approx0.841471$ while $q_4(1)=\tfrac14\sin 1\approx0.210368$.
These predictions differ, so by EL-D2 the objects are not $\sim_P$-equivalent.

For the consequence: a well-defined assignment $\Phi$ from $\cong$-classes to
$\sim_P$-classes would give
$[\mathcal{R}_1]_{\sim_P}=\Phi([\mathcal{R}_1]_{\cong})
=\Phi([\mathcal{R}_2]_{\cong})=[\mathcal{R}_2]_{\sim_P}$, contradicting the
previous paragraph.
\end{proof}

\begin{interpretation}[What blocks descent, S-I2]
\label{int:blocks}
The map $F$ is a perfectly good symplectomorphism; what it fails to do is carry
the calibration. Under $F$ the numeral read off the port is multiplied by
$\alpha$, so ``$1$ metre'' in $\mathcal{R}_1$ becomes ``$\alpha$ metres'' in
$\mathcal{R}_2$ unless the metre itself is transported. Once it is, EL-D3 holds
and EL-P2 restores $\sim_P$. S-N2 is therefore not in tension with EL-P2: it
identifies which hypothesis of EL-P2 is doing the work.
\end{interpretation}

\begin{corollary}[Ports are load-bearing, S-C2]
\label{cor:ports}
The interface component $I$ of Definition~\ref{def:object} cannot be omitted
without loss: two objects agreeing on $K$ and $D$ up to isomorphism and
disagreeing on $I$ can have different protocol-equivalence classes.
\end{corollary}

\begin{proof}
The two objects of S-N2 agree on $(K,D)$ up to $\cong$. They differ in which
function of the state is designated as the calibrated port, since the port of
$\mathcal{R}_2$ pulled back along $F$ is $\alpha^{-1}$ times the port of
$\mathcal{R}_1$. S-N2 shows their $\sim_P$-classes differ.
\end{proof}

\section{Limits need not commute with composition}
\label{sec:limits}

EL proves that the Newtonian and free-particle limits hold only after a domain
is fixed (EL-P3, EL-P4). CS proves that an interaction is not determined by the
factors (CS-P4). Together these obstruct a natural expectation: that one may
take a limit factor by factor and then compose.

\begin{newresult}[Non-commutation of limit and composition, S-N3]
\label{res:limits}
Let $\Sigma_\omega(\kappa)$ denote the composite of two identical
one-degree-of-freedom oscillators of mass $m$ and frequency $\omega$ with
reciprocal coupling $\kappa$, that is
$H_{\omega,\kappa}=\frac{p_1^2+p_2^2}{2m}
+\frac{m\omega^2(q_1^2+q_2^2)}{2}+\kappa q_1q_2$. Let
$\Sigma^{(i)}_\omega$ denote the isolated factors. For every $\kappa\neq0$,
\[
 \lim_{\omega\to0}\Sigma_\omega(\kappa)\;\neq\;
 \Big(\lim_{\omega\to0}\Sigma^{(1)}_\omega\Big)\otimes
 \Big(\lim_{\omega\to0}\Sigma^{(2)}_\omega\Big),
\]
where $\otimes$ is the composition of Definition~\ref{def:object} supplied with
the interaction data that the factors determine, namely none. Moreover the two
sides differ qualitatively and not merely perturbatively: with initial data
$(q_1,p_1,q_2,p_2)=(0,0,1,0)$ and $\kappa=m=1$, the left side gives
$q_1(t)=\tfrac12(\cos t-\cosh t)$ and the right side gives $q_1(t)\equiv0$.
\end{newresult}

\begin{proof}
The left side has Hamiltonian $\frac{p_1^2+p_2^2}{2m}+\kappa q_1q_2$, obtained
by setting $\omega=0$ in $H_{\omega,\kappa}$. Each isolated factor has
$\omega\to0$ limit $\frac{p_i^2}{2m}$, a free particle. By CS-P4 the pair of
factors does not determine an interaction term, so the composition on the right
supplies none and its Hamiltonian is $\frac{p_1^2+p_2^2}{2m}$. The two
Hamiltonians differ by $\kappa q_1q_2$, whose differential
$\kappa(q_2\,dq_1+q_1\,dq_2)$ is nonzero away from $q_1=q_2=0$; by nondegeneracy
of $\omega_\oplus$ the generated fields differ there.

For the explicit trajectories, take $\kappa=m=1$. The left side gives
$\ddot q_1=-q_2$, $\ddot q_2=-q_1$. In the variables $u=q_1+q_2$ and
$v=q_1-q_2$ these decouple to $\ddot u=-u$ and $\ddot v=+v$. The stated initial
data give $u(0)=1,\dot u(0)=0$ and $v(0)=-1,\dot v(0)=0$, hence $u=\cos t$ and
$v=-\cosh t$, so $q_1=\tfrac12(u+v)=\tfrac12(\cos t-\cosh t)$. The right side
gives $\ddot q_1=0$ with $q_1(0)=\dot q_1(0)=0$, hence $q_1\equiv0$. At $t=1$
the values are $-0.5013891\ldots$ and $0$.
\end{proof}

\begin{empirical}[Numerical cross-check, S-E2]
\label{emp:numeric}
Independent Runge--Kutta integration of the left-hand system in
\texttt{Core.hs}, using $2\times10^{4}$ steps of a fourth-order scheme on
$[0,1]$, returns $q_1(1)=-0.501389$, agreeing with the closed form of S-N3 to
the six digits printed. This is an agreement between two computations performed
in this round; it is not an observation of any physical system.
\end{empirical}

\begin{failuremode}[Why the diagram fails]
The failure is not an analytic subtlety about the order of limits. It is the
interface layer asserting itself: $\kappa$ lives in $I$, not in either factor's
$K$ or $D$, so no factorwise operation can see it. Whenever a limiting relation
is claimed compatible with composition, the claim must say what happens to the
interaction data, and that is extra input.
\end{failuremode}

\begin{counterexample}[The instability is not an artifact of taking $\omega=0$]
For small $\omega>0$ the same normal-mode computation gives
$\ddot v=(\kappa/m-\omega^2)v$, exponentially unstable whenever
$\kappa>m\omega^2$. The limit does not create the instability; it removes the
last restoring term that could have suppressed it. Reporting ``the $\omega\to0$
limit of a coupled pair is a pair of free particles'' discards a growing mode.
\end{counterexample}

\section{What the formal drafts establish, and what they do not}
\label{sec:formal}

Two source drafts accompany this paper. Both were executed in this round; the
results are reported as process facts, not as physical facts.

\begin{empirical}[Toolchain record, S-E3]
\label{emp:toolchain}
\texttt{formal/drafts/Proofs.lean} was checked by \texttt{lean} version
$4.33.0$ (arm64 Darwin) with no Mathlib dependency, exiting $0$ with no errors
and no warnings; it contains no \texttt{sorry}, and \texttt{\#print axioms} on
each named theorem reports dependence on at most \texttt{propext} and
\texttt{Quot.sound}. \texttt{formal/drafts/Core.hs} type-checks under
\texttt{ghc} $8.4.2$ with \texttt{-Wall} and no warnings, and runs to completion
under \texttt{runghc} with $21$ of $21$ declared witnesses passing. Native
linking of \texttt{Core.hs} was \emph{not} achieved in this environment: the
installed GHC is an x86\_64 build invoked on an arm64 host, and its assembler
stage fails on generated x86 assembly. That failure is a property of the local
installation, not of the source; type-checking and interpreted execution both
succeed.
\end{empirical}

The Lean development formalizes an abstraction, and the abstraction is lossy in
ways that must be stated. It contains typed systems and processes with an
observable pullback; the category laws of CS-P1; product composition with
functoriality and an involutive symmetry; nondegeneracy as an explicit
bijection, with the resulting factorization statement corresponding to CS-T1;
the separable/non-separable distinction with both classes inhabited,
corresponding to CS-P4; the scalar coupling trichotomy of S-N1 and the split
obstruction of S-P2 over $\mathbb{Z}$; and a protocol layer comprising
observational equivalence with its equivalence-relation laws, the prediction-map
characterization corresponding to EL-N1, the transported-interpretation bridge
corresponding to EL-P2, a witness that a bare bijection does not suffice, the
blind-protocol degeneracy, and the falsification results corresponding to EL-T1
and EL-C2.

It does not contain smooth manifolds, differential forms, the closedness
condition $d\omega=0$, de Rham cohomology, flows, or any analytic limit. Every
Lean statement whose informal counterpart lives in differential geometry is an
abstraction in which the geometric content has been replaced by a declared
interface; in particular the closedness half of CS-L1 is absent, HR-T1 is
absent, and only the \emph{nondegeneracy} half of the product-symplectic lemma
is represented. No theorem in this synthesis should be described as
machine-checked in its geometric form. The Haskell draft contains no proofs at
all: it supplies decision procedures for reciprocity and for the coupling class
of S-N1, a witness generator for S-P2, a Runge--Kutta comparison of composite
and factorized flows, the transport and port computations of S-N2, the
integration underlying S-E2, and executable versions of EL's counterexamples on
approximate transitivity, finite sampling, and bounded-domain limits. Each is a
finite computation reported alongside its declared tolerance, and a passing
witness establishes nothing universal.

\section{Preserved disagreements, gaps, and failed constructions}
\label{sec:disagreements}

\begin{disagreement}[Two incompatible proposals for empirical composition]
CS-C2 conjectures that empirical composition should be organized by \emph{lax
symmetric-monoidal realization maps} into an operational category. EL-C3
conjectures instead that it should be organized by \emph{protocol-composition
maps} under which predictions factorize \emph{up to a quantified tolerance} when
the systems are operationally isolated to that tolerance. These are not the same
proposal: the first is exact and structural and presupposes a target category
whose existence CS explicitly does not establish; the second is approximate and
metric and presupposes a tolerance calculus that EL does not construct. This
synthesis does not adjudicate between them, does not show them equivalent, and
notes that the non-transitivity of $\varepsilon$-equivalence is an obstacle
specific to the second.
\end{disagreement}

\begin{disagreement}[The numerical coefficient in HR-P1]
HR-P1 asserts $d\omega(X_f,X_g,X_h)=2\operatorname{Jac}(f,g,h)$. The
\emph{qualitative} biconditional it supports, that a nondegenerate two-form is
closed if and only if its induced bracket satisfies the Jacobi identity, is
standard and is what Sections~\ref{sec:trichotomy} and \ref{sec:limits} use. The
coefficient $2$ and the overall sign, however, depend jointly on the sign
convention in $\iota_{X_H}\omega=\pm dH$ and on the normalization convention for
the exterior derivative of a two-form. This synthesis did not re-derive the
coefficient and does not endorse it. Nothing proved here depends on it. HR's
accompanying counterexample \emph{was} re-derived: for the bracket on $\R^3$
with $\{x_1,x_2\}=x_1$, $\{x_2,x_3\}=x_3$, $\{x_3,x_1\}=x_1$, the Jacobiator on
$(x_1,x_2,x_3)$ is $\{x_1,x_3\}+\{x_3,x_1\}+\{x_1,x_2\}=x_1$, confirming that
Leibniz does not imply Jacobi.
\end{disagreement}

\begin{disagreement}[Scope of the review record]
Empirical statement S-E1 records that the supplied review returned
\texttt{ACCEPT} with confidence $1$ and uniform per-claim explanations. This
synthesis treats that record as a process fact and not as verification, and
consequently re-derived the claims marked ``re-derived'' in
Table~\ref{tab:inherited} rather than relying on it. Claims marked ``inherited''
in that table therefore rest on the source papers' own arguments and on nothing
else in this round.
\end{disagreement}

\begin{failuremode}[Bibliographic verification not performed]
The prior artifact records that bibliographic records not independently
rechecked retain the verification status documented by the source workers, and
that the pilot's source-candidate list marks all four of its entries with the
status ``candidate requires agent verification''. No network access was used in
this round, so no bibliographic identity was confirmed here. The bibliography
below is carried forward at the sources' declared status.
\end{failuremode}

\begin{failuremode}[Constructions attempted and not completed]
Three constructions were attempted in drafting this synthesis and are recorded
as incomplete rather than omitted. First, a functor from realized theory objects
to protocol quotients: S-N2 shows that any such functor must take the empirical
layer as an argument, and no candidate with the right universal property was
found. Second, an extension of the trichotomy S-N1 to nonlinear interconnections:
the contraction of Lemma~\ref{lem:contraction} is specific to the linear family
\eqref{eq:Y}, and no general criterion was obtained. Third, a Lean formalization
of the closedness half of CS-L1: this requires differential forms and was not
attempted without Mathlib.
\end{failuremode}

\begin{conjecture}[Interface-indexed composition, S-C3]
\label{conj:interface}
The negative results S-N1, S-N2, and S-N3 point to a common positive shape,
which is here left unproved: that composition of realized theory objects is
functorial once it is indexed by interface data, that is, that fixing the port
list $I$ and the interaction pattern turns $\otimes$ into an operation for which
protocol quotients, isomorphism classes, and limits all behave compatibly. No
category has been exhibited in which this is true, no coherence conditions have
been checked, and the tolerance calculus that S-N3 would require is absent.
\end{conjecture}

\section{Falsification exposure of this synthesis}
\label{sec:falsification}

Per EL-D5, a claim is testable only if what would count against it is stated in
advance.

\begin{enumerate}[leftmargin=2.2em]
\item S-N1 is refuted by exhibiting $(C_{12},C_{21})$ with $C_{21}=C_{12}^{T}$,
      $C_{12}\neq0$, whose total field nevertheless splits, or by exhibiting a
      non-reciprocal pair whose field preserves $\omega_\oplus$. Either would
      contradict Lemma~\ref{lem:contraction}, which is a finite computation and
      can be checked directly.
\item S-P2 is refuted by exhibiting $a_1,a_2\neq0$ and a unidirectional
      $(c_{12},c_{21})$ with $A^{T}\Omega_0+\Omega_0A=0$. The matrix $M$
      displayed in the proof makes this a four-by-four inspection.
\item S-N2 is refuted by showing that the two oscillator objects are in fact
      $\sim_P$-equivalent under the declared protocol, which would require
      $\sin 1=\tfrac14\sin 1$; or by showing that the declared empirical layer
      is inadmissible under EL-D1, which would change the claim's scope rather
      than its truth.
\item S-N3 is refuted by showing that the composition of
      Definition~\ref{def:object} does determine an interaction from the factors
      alone, contradicting CS-P4; or that $\tfrac12(\cos t-\cosh t)\equiv0$.
\item S-E2 and S-E3 are refuted by re-running the recorded commands and
      obtaining different output; both name the tools, versions, and exit
      conditions required to do so.
\end{enumerate}

None of these is an empirical test of a physical claim, because this paper
advances no physical claim. That is itself a limitation: the synthesis operates
within the mathematical and process layers, and its contact with the empirical
layer is confined to saying what would have to be supplied for such contact to
exist.

\section{Conclusion}

Three separations survive the synthesis. Composability splits into a trichotomy
in which being Hamiltonian and being factorized are distinct conditions, with a
nonempty class between them. Mathematical isomorphism does not descend to
protocol-relative observational equivalence, so the carrier of empirical content
is the pair (mathematical layer, empirical layer) and never the mathematical
layer alone. Limiting relations do not commute with composition, because the
interaction data lives in an interface layer that factorwise operations cannot
see. Each separation supports the negative half of the founding thesis, that
``compositional'' and ``empirically realized'' are not interchangeable
predicates; none supports its positive half. Whether empirically realized
compositional structure admits a single theory with a universal property remains
conjectural, and Conjecture~\ref{conj:interface} states what would have to be
built.

\paragraph{Note on the bibliography.} Every entry below is carried forward from
CS, HR, or EL at the verification status those papers declared. No entry was
independently rechecked in this round; see the failure mode recorded in
Section~\ref{sec:disagreements}.

{\footnotesize
\setlength{\itemsep}{0pt}
\begin{thebibliography}{99}
\setlength{\itemsep}{0pt}
\setlength{\parskip}{0pt}
\bibitem{AbrahamMarsden1978} R. Abraham and J. E. Marsden, \emph{Foundations of
Mechanics}, 2nd ed., AMS Chelsea, 2008.
\bibitem{MarsdenRatiu1999} J. E. Marsden and T. S. Ratiu, \emph{Introduction to
Mechanics and Symmetry}, 2nd ed., Springer, 1999.
\bibitem{McDuffSalamon2017} D. McDuff and D. Salamon, \emph{Introduction to
Symplectic Topology}, 3rd ed., Oxford University Press, 2017.
\bibitem{MacLane1998} S. Mac Lane, \emph{Categories for the Working
Mathematician}, 2nd ed., Springer, 1998.
\bibitem{BaezStay2011} J. Baez and M. Stay, ``Physics, Topology, Logic and
Computation: A Rosetta Stone,'' 2011, \url{https://arxiv.org/abs/0903.0340}.
\bibitem{vanderSchaft2014} A. van der Schaft and D. Jeltsema, ``Port-Hamiltonian
Systems Theory: An Introductory Overview,'' \emph{Foundations and Trends in
Systems and Control} 1 (2014), 173--378.
\bibitem{Onsager1931} L. Onsager, ``Reciprocal Relations in Irreversible
Processes I,'' \emph{Physical Review} 37 (1931), 405--426.
\bibitem{Weatherall2019} J. O. Weatherall, ``Theoretical Equivalence in Physics,''
Parts 1 and 2, \emph{Philosophy Compass} 14 (2019), e12592 and e12591.
\bibitem{Coffey2014} K. Coffey, ``Theoretical Equivalence as Interpretative
Equivalence,'' \emph{British Journal for the Philosophy of Science} 65 (2014),
821--844.
\bibitem{Butterfield2011} J. Butterfield, ``Less is Different: Emergence and
Reduction Reconciled,'' \emph{Foundations of Physics} 41 (2011), 1065--1135.
\end{thebibliography}}

\end{document}
